% Generated draft: requires mathematical and source review.
\documentclass[11pt]{article}
\usepackage[margin=25mm]{geometry}
\usepackage{fontspec}
\setmainfont{Latin Modern Roman}
\usepackage{amsmath,amssymb,mathtools,unicode-math}
\setmathfont{Latin Modern Math}
\usepackage{enumitem,needspace,xurl,hyperref}
\setlength{\parindent}{0pt}
\setlength{\parskip}{5pt}
\setlength{\emergencystretch}{3em}
\newcommand{\N}{\mathbb{N}}
\newcommand{\Z}{\mathbb{Z}}
\newcommand{\Q}{\mathbb{Q}}
\newcommand{\R}{\mathbb{R}}
\newcommand{\C}{\mathbb{C}}
\newcommand{\F}{\mathbb{F}}
\newcommand{\A}{\mathbb{A}}
\newcommand{\PP}{\mathbb{P}}
\newcommand{\E}{\mathbb{E}}
\newcommand{\eps}{\varepsilon}
\newcommand{\dd}{\,\mathrm{d}}
\begin{document}
\section*{UnsolvedMath ID: 30004522}
\textbf{Source problem code:} catalogue entry
\subsection{Definitions and mathematical statement}
The supplied record combines five questions. The first has a standard precise formulation; the others require the qualifications recorded below.

\begin{enumerate}
\item Let \(P\subset\R^4\) be a full-dimensional convex polytope, let \(V(P)\) be its vertex set, and put \(n=|V(P)|\). The polytope is simplicial if each facet is a tetrahedron. Its graph \(G(P)\) has vertex set \(V(P)\), with two vertices adjacent exactly when their segment is an edge of \(P\). Assume that
\[
\forall T\subseteq V(P),\quad |T|=3\ \Longrightarrow\ \exists\,u,v\in T,\quad u\ne v,\quad \{u,v\}\in E(G(P)).
\]
Must there exist \(K\subseteq V(P)\) such that
\[
|K|\ge\left\lceil\frac n{10}\right\rceil,
\qquad
\forall\,u,v\in K,\quad u\ne v\ \Longrightarrow\ \{u,v\}\in E(G(P))?
\]
This question is quantified over every simplicial convex \(4\)-polytope satisfying the displayed hypothesis.

\item Write \(S^2=\{x\in\R^3:\|x\|=1\}\). For \(u\in S^2\) and \(-1\le t\le1\), define the closed and open spherical caps
\[
C_{\ge}(u,t)=\{x\in S^2:\langle u,x\rangle\ge t\},
\qquad
C_{>}(u,t)=\{x\in S^2:\langle u,x\rangle>t\}.
\]
The source asks whether, for every pair \(\mu_1,\mu_2\) of continuous probability measures on \(S^2\) and every \(\alpha\in(0,1/2)\), there is a spherical cap \(C\) satisfying
\[
\mu_1(C)=\mu_2(C)=\alpha.
\]
The source does not define ``continuous probability measure'' or choose an open or closed cap convention. Thus this is a qualified question, not a fully specified assertion: the admissible measure class and cap convention must be supplied. In particular, absence of point masses does not by itself ensure that cap boundaries have measure zero. No choice between atomless measures, measures absolutely continuous with respect to surface area, or measures with continuous densities is made here.

\item Under a usual finite affine interpretation, a pseudoline is a properly embedded copy of \(\R\) in \(\R^2\); a simple pseudoline arrangement is a finite family in which each pair crosses exactly once and no three pseudolines meet. Color its pseudolines red or blue. A face is a connected component of the complement of the union of the pseudolines. A bounded triangular face is a bounded face with three boundary edges. Interpreting incidence with both colors to mean that its boundary edges lie on pseudolines of both colors, the supplied question asks whether every such arrangement has such a face.

These definitions give a conventional interpretation, but the source does not specify the affine or projective setting, finiteness, minimum size, or whether both colors must occur. The unrestricted formulation includes monochromatic arrangements and arrangements with fewer than three pseudolines, so it cannot be an affirmative universal assertion as written. Even requiring both colors leaves two-pseudoline arrangements without a triangular face. The intended hypotheses must be recovered before identifying this question with the theorem mentioned in the supplied literature assessment.

\item For a hypergraph \(H=(V,E)\), where \(E\) is a family of subsets of \(V\), the source asks whether
\[
m_2^*(H)=3\quad\Longrightarrow\quad m_3^*(H)<\infty.
\]
Here the arguments \(H\) merely display the dependence suggested by the source phrase ``for a hypergraph \(H\).'' Neither \(m_2^*\) nor \(m_3^*\) is defined in the record. The class of hypergraphs, finiteness assumptions, and conventions for these parameters are also absent. Consequently the displayed implication is only a transcription of the intended question, not a complete mathematical formulation.

\item Let \(X\subset\R^2\) be a set of \(n\) distinct points in general position, interpreted here in the usual planar sense that no three are collinear. Once the intended definition of a partial triangulation is fixed, let \(T_{\mathrm{part}}(X)\) denote its number of partial triangulations. The source asks for
\[
\min_{\substack{X\subset\R^2,\ |X|=n\\\text{no three points of }X\text{ collinear}}}
T_{\mathrm{part}}(X)
\]
and asks whether this minimum equals
\[
C_{n-2}=\frac{1}{n-1}\binom{2n-4}{n-2}.
\]
The formula is stated for integer \(n\ge2\), its usual numerical range here; the intended geometric range of \(n\) is not supplied. Crucially, ``partial triangulation'' is undefined: the record does not say whether all points must be used, whether the convex hull must be covered, which boundary edges are required, or whether an arbitrary noncrossing subset of triangulation edges is counted. The counting convention and any additional general-position assumptions require confirmation before this minimum defines the intended problem.
\end{enumerate}
\subsection{Short English statement}
This combined record asks about large cliques in simplicial \(4\)-polytope graphs, equal-measure spherical caps, triangles with both colors in pseudoline arrangements, an implication between undefined hypergraph parameters, and the minimum number of partial triangulations. Several hypotheses and definitions are missing. The supplied dated curation labels the combined record open and reports a proof of the pseudoline assertion, but the precise scope and literature claims require independent review.
\subsection{Sources}
\begin{itemize}
\item Local registry: \texttt{data/unsolvedmath/problems.json}, permanent UnsolvedMath ID \(30004522\), title ``Extremal Questions in Discrete Geometry and Topological Combinatorics,'' exact source problem code \texttt{catalogue entry}. The supplied record belongs to the Oberwolfach Reports Open Problems collection and gives proposed year \(2020\); no individual proposer is supplied.
\item Supplied report citation: \emph{Discrete Geometry (2021)}, Oberwolfach Reports, DOI \texttt{10.4171/owr/2020/30}, pp.\ 1469--1528. \url{https://doi.org/10.4171/owr/2020/30}. The background describes this as a workshop report containing open-problem discussions. Its contents were not independently consulted for this draft.
\item Supplied literature source: Lauff--Radtke, \emph{On Triangles in Colored Pseudoline Arrangements}. \url{https://arxiv.org/abs/2601.20574}.
\item Supplied literature source: \emph{Oberwolfach Reports 2020/30}. \url{https://ems.press/content/serial-article-files/46867}.
\item Literature assessment checked at \texttt{2026-08-22}: this is supplied dated curation, not independently verified. It labels the combined record open, states that the four-polytope clique and spherical-cap questions remain open, and reports that a \(2026\) preprint proves the bichromatic-triangle assertion. It does not supply individual status assessments for the hypergraph or partial-triangulation questions. The background expressly requests independent verification of time-sensitive status and cited claims.
\item The registry supplies the statement-status label \texttt{corrected\_verified}, qualified by the explanation that only formula delimiters, notation, grammar, and typography were repaired while mathematical content and scope were preserved. This label is not treated here as independent mathematical verification.
\end{itemize}
\end{document}
